%%%%%%%%%%%%%%%%
%%% gen 41 %%%
%%%%%%%%%%%%%%%%

\Chapter{41}

\Verse{1}
{
	\hE{וַיְהִי מִקֵּץ שְׁנָתַיִם יָמִים וּפַרְעֹה חֹלֵם וְהִנֵּה עֹמֵד עַל הַיְאֹר׃}
}
{
	\eE{
		And it was at the end of two years’ time, and Pharaoh was
		dreaming: and here he was standing by the Nile.
	}
}
{
	(Narrator: Two years later.)

	The king's having a dream.

	He's out by the river.
}

\Verse{2}
{
	\hE{וְהִנֵּה מִן הַיְאֹר עֹלֹת שֶׁבַע פָּרוֹת יְפוֹת מַרְאֶה וּבְרִיאֹת בָּשָׂר וַתִּרְעֶינָה בָּאָחוּ׃}
}
{
	\eE{
		And here, coming up from the Nile, were seven cows,
		beautiful-looking and fat-fleshed, and they fed in the
		reeds.
	}
}
{
	He sees a bible number of cows.
	
	Seven this time.
	
	Fat as a cow, they were.
}

\Verse{3}
{
	\hE{וְהִנֵּה שֶׁבַע פָּרוֹת אֲחֵרוֹת עֹלוֹת אַחֲרֵיהֶן מִן הַיְאֹר רָעוֹת מַרְאֶה וְדַקּוֹת בָּשָׂר וַתַּעֲמֹדְנָה אֵצֶל הַפָּרוֹת עַל שְׂפַת הַיְאֹר׃}
}
{
	\eE{
		And here were seven other cows coming up after them from
		the Nile, bad-looking and thin-fleshed, and they stood by
		the cows that were by the bank of the Nile.
	}
}
{
	Then there's seven skinny ones.

	They come stand by the good cows, uninvited.
}

\Verse{4}
{
	\hE{וַתֹּאכַלְנָה הַפָּרוֹת רָעוֹת הַמַּרְאֶה וְדַקֹּת הַבָּשָׂר אֵת שֶׁבַע הַפָּרוֹת יְפֹת הַמַּרְאֶה וְהַבְּרִיאֹת וַיִּיקַץ פַּרְעֹה׃}
}
{
	\eE{
		And the bad-looking and thin-fleshed cows ate the seven
		beautiful-looking and fat cows—and Pharaoh woke up.
	}
}
{
	Then the skinny ones eat the fat ones.

	Cuz obviously.

	They were hungry.
	
	The king wakes up with a start.
}

\Verse{5}
{
	\hE{וַיִּישָׁן וַיַּחֲלֹם שֵׁנִית וְהִנֵּה שֶׁבַע שִׁבֳּלִים עֹלוֹת בְּקָנֶה אֶחָד בְּרִיאוֹת וְטֹבוֹת׃}
}
{
	\eE{
		And he slept and dreamed a second time, and here were
		seven ears of grain coming up on one stalk, fat and good.
	}
}
{
	Back to bed.
	
	Second dream.
	
	There's a bible number of things again.

	Seven again.

	This time it's a bread plant.
	
	With seven breadsticks.
}

\Verse{6}
{
	\hE{וְהִנֵּה שֶׁבַע שִׁבֳּלִים דַּקּוֹת וּשְׁדוּפֹת קָדִים צֹמְחוֹת אַחֲרֵיהֶן׃}
}
{
	\eE{
		And here were seven ears of grain, thin and scorched by
		the east wind, growing after them.
	}
}
{
	Then there's seven other breadsticks.

	They're burnt.
}

\Verse{7}
{
	\hE{וַתִּבְלַעְנָה הַשִּׁבֳּלִים הַדַּקּוֹת אֵת שֶׁבַע הַשִּׁבֳּלִים הַבְּרִיאוֹת וְהַמְּלֵאוֹת וַיִּיקַץ פַּרְעֹה וְהִנֵּה חֲלוֹם׃}
}
{
	\eE{
		And the thin ears of grain swallowed the seven fat and
		full ears of grain—and Pharaoh woke up, and here it was a
		dream.
	}
}
{
	The bad breadsticks eat the good ones.

	Pharaoh wakes up again in a cold sweat.

	And here it was a dream.
}

\Verse{8}
{
	\hE{וַיְהִי בַבֹּקֶר וַתִּפָּעֶם רוּחוֹ וַיִּשְׁלַח וַיִּקְרָא אֶת כּל חַרְטֻמֵּי מִצְרַיִם וְאֶת כּל חֲכָמֶיהָ וַיְסַפֵּר פַּרְעֹה לָהֶם אֶת חֲלֹמוֹ וְאֵין פּוֹתֵר אוֹתָם לְפַרְעֹה׃}
}
{
	\eE{
		And it was in the morning, and his spirit was moved, and
		he sent and called all of Egypt’s magicians and all its
		wise men. And Pharaoh told them his dream, and there was
		no one who could tell the meaning of them to Pharaoh.
	}
}
{
	Pharaoh's scared.

	He calls \textsl{all the magicians in the country}.

	Also all the wise guys.

	Nobody knows what it means, and no one's willing to guess.
}

\Verse{9}
{
	\hE{וַיְדַבֵּר שַׂר הַמַּשְׁקִים אֶת פַּרְעֹה לֵאמֹר אֶת חֲטָאַי אֲנִי מַזְכִּיר הַיּוֹם׃}
}
{
	\eE{
		And the chief of the drink-stewards spoke to Pharaoh,
		saying, “I’m recalling my sins today:
	}
}
{
	The wi(n/s)e guy is like ``Hey remember when I was bad?''
}

\Verse{10}
{
	\hE{פַּרְעֹה קָצַף עַל עֲבָדָיו וַיִּתֵּן אֹתִי בְּמִשְׁמַר בֵּית שַׂר הַטַּבָּחִים אֹתִי וְאֵת שַׂר הָאֹפִים׃}
}
{
	\eE{
		Pharaoh had been angry at his servants, and he put me
		under watch at the house of the chief of the guards, me
		and the chief of the bakers.
	}
}
{
	``And you put me in jail with that (d+br)ead guy?''
}

\Verse{11}
{
	\hE{וַנַּחַלְמָה חֲלוֹם בְּלַיְלָה אֶחָד אֲנִי וָהוּא אִישׁ כְּפִתְרוֹן חֲלֹמוֹ חָלָמְנוּ׃}
}
{
	\eE{
		And we had a dream in one night, I and he; we each dreamed
		with his own dream’s meaning.
	}
}
{
	``Well back then, \textsl{we} had a dream!''\fC{
		There was a brief conversation here clarifying
		some issues, which we've omitted for brevity.
	}
}

\Verse{12}
{
	\hE{וְשָׁם אִתָּנוּ נַעַר עִבְרִי עֶבֶד לְשַׂר הַטַּבָּחִים וַנְּסַפֶּר לוֹ וַיִּפְתּר לָנוּ אֶת חֲלֹמֹתֵינוּ אִישׁ כַּחֲלֹמוֹ פָּתָר׃}
}
{
	\eE{
		And there, with us, was a Hebrew boy, a slave of the chief
		of the guards; and we told him, and he told us the meaning
		of our dreams. He told the meaning according to each one’s
		dream.
	}
}
{
	``There was this Jew too.
	
	One of Potiphar's guys.

	He told us what the dreams meant.

	And now the other guy is wicked (d+br)ead just like he said.''
}

\Verse{13}
{
	\hE{וַיְהִי כַּאֲשֶׁר פָּתַר לָנוּ כֵּן הָיָה אֹתִי הֵשִׁיב עַל כַּנִּי וְאֹתוֹ תָלָה׃}
}
{
	\eE{
		And it was—as he told the meaning for us, so it was: me he
		put back at my station, and him he hanged.”
	}
}
{
	``I got my job back, and he's toast.''
}

\Verse{14}
{
	\hE{וַיִּשְׁלַח פַּרְעֹה וַיִּקְרָא אֶת יוֹסֵף וַיְרִיצֻהוּ מִן הַבּוֹר וַיְגַלַּח וַיְחַלֵּף שִׂמְלֹתָיו וַיָּבֹא אֶל פַּרְעֹה׃}
}
{
	\eE{
		And Pharaoh sent and called Joseph, and they rushed him
		from the pit, and he shaved and changed his clothing and
		came to Pharaoh.
	}
}
{
	So Pharaoh makes some calls and they pull Joseph out of the pit.\fA{
		The pit is the jail. Jails were different back then.
	}

	They give him a makeover\fA{He's been in the pit for a couple years.}
	and bring him to the king.
}

\Verse{15}
{
	\hE{וַיֹּאמֶר פַּרְעֹה אֶל יוֹסֵף חֲלוֹם חָלַמְתִּי וּפֹתֵר אֵין אֹתוֹ וַאֲנִי שָׁמַעְתִּי עָלֶיךָ לֵאמֹר תִּשְׁמַע חֲלוֹם לִפְתֹּר אֹתוֹ׃}
}
{
	\eE{
		And Pharaoh said to Joseph, “I’ve had a dream, and there’s
		no one who can tell the meaning of it, and I’ve heard
		about you, saying you hear a dream so as to tell the
		meaning of it.”
	}
}
{
	The Pharaoh's like ``I've had a dream. I hear you're dreamy.''
}

\Verse{16}
{
	\hE{וַיַּעַן יוֹסֵף אֶת פַּרְעֹה לֵאמֹר בִּלְעָדָי אֱלֹהִים יַעֲנֶה אֶת שְׁלוֹם פַּרְעֹה׃}
}
{
	\eE{
		And Joseph answered Pharaoh, saying, “Not I. {\GodInit} will
		answer regarding Pharaoh’s well-being.”
	}
}
{
	Joseph's give him the ``It's not me it's {\God}, but I
	can talk to him for you'' line again.\footnote{
		A well known trick between monarchs and their viziers,
		see also Dan 2:27–30 and Ala 1:05:22; timestamps
		vary between manuscripts.
	}
}

\Verse{17}
{
	\hE{וַיְדַבֵּר פַּרְעֹה אֶל יוֹסֵף בַּחֲלֹמִי הִנְנִי עֹמֵד עַל שְׂפַת הַיְאֹר׃}
}
{
	\eE{
		And Pharaoh spoke to Joseph, “In my dream, here I was
		standing by the bank of the Nile.
	}
}
{
	He's like ``I was by the river.''
}

\Verse{18}
{
	\hE{וְהִנֵּה מִן הַיְאֹר עֹלֹת שֶׁבַע פָּרוֹת בְּרִיאוֹת בָּשָׂר וִיפֹת תֹּאַר וַתִּרְעֶינָה בָּאָחוּ׃}
}
{
	\eE{
		And here, coming up from the Nile, were seven cows,
		fat-fleshed and beautiful-figured, and they fed in the
		reeds.
	}
}
{
	He explains the cows.
}

\Verse{19}
{
	\hE{וְהִנֵּה שֶׁבַע פָּרוֹת אֲחֵרוֹת עֹלוֹת אַחֲרֵיהֶן דַּלּוֹת וְרָעוֹת תֹּאַר מְאֹד וְרַקּוֹת בָּשָׂר לֹא רָאִיתִי כָהֵנָּה בְּכל אֶרֶץ מִצְרַיִם לָרֹעַ׃}
}
{
	\eE{
		And here were seven other cows coming up after them, weak
		and very bad-figured and scrawny-fleshed. I haven’t seen
		any like these in all the land of Egypt for bad!
	}
}
{
	And the other cows.

	\aB{
	Bible note: in the thing about the cows and the thing
	about the breadsticks, it appears that the scribal
	bible people who copied these books by hand for
	all those years before Gutenberg repeatedly confused
	the letters D and R in Hebrew, since they looked
	like \paleo{ד} and \paleo{ר} in paleo-Hebrew,
	\ara{𐡃} and \ara{𐡓} in Aramaic, and \heb{ד} and \heb{ר}
	in the modern Aramaic-derived square script that we call
	Hebrew for some reason.
	}


	\aC{
	Consult the following chart for examples of the Dalet/Resh
	phenomenon in this and the surrounding chapters.

		\Table{|p{0.10\linewidth}|p{0.20\linewidth}|p{0.20\linewidth}|p{0.40\linewidth}|}{
			\hline
			\textbf{Verse}
				& \textbf{Hebrew}
				& \textbf{Consonants}
				& \textbf{Meaning} \\
			\hline
			41:3
				& \heb{דַקּוֹת בָּשָׂר}
				& D-Q
				& ``thin-fleshed'' cows \\
			\hline
			41:4
				& \heb{דַקֹּת הַבָּשָׂר}
				& D-Q
				& ``thin-fleshed'' cows \\
			\hline
			41:6
				& \heb{דַּקּוֹת}
				& D-Q
				& ``thin'' ears of grain \\
			\hline
			41:7
				& \heb{הַדַּקּוֹת}
				& D-Q
				& ``the thin'' ears of grain \\
			\hline
			41:19
				& \heb{רַקּוֹת בָּשָׂר}
				& R-Q
				& ``scrawny-fleshed'' cows \\
			\hline
			41:20
				& \heb{הָרַקּוֹת}
				& R-Q
				& ``the scrawny'' cows \\
			\hline
			41:23
				& \heb{דַּקּוֹת}
				& D-Q
				& ``thin'' ears of grain \\
			\hline
			41:24
				& \heb{הַדַּקֹּת}
				& D-Q
				& ``the thin'' ears of grain \\
			\hline
			41:27
				& \heb{הָרַקּוֹת}
				& R-Q
				& ``the scrawny'' cows \\
			\hline
			41:27
				& \heb{הָרֵקוֹת}
				& R-Q
				& ``the empty/scrawny'' ears of grain \\
			\hline
		}

	}
}

\Verse{20}
{
	\hE{וַתֹּאכַלְנָה הַפָּרוֹת הָרַקּוֹת וְהָרָעוֹת אֵת שֶׁבַע הַפָּרוֹת הָרִאשֹׁנוֹת הַבְּרִיאֹת׃}
}
{
	\eE{
		And the cows that were scrawny and bad ate the first seven
		cows that were fat,
	}
}
{
	The hungry cows ate the tasty cows.
}

\Verse{21}
{
	\hE{וַתָּבֹאנָה אֶל קִרְבֶּנָה וְלֹא נוֹדַע כִּי בָאוּ אֶל קִרְבֶּנָה וּמַרְאֵיהֶן רַע כַּאֲשֶׁר בַּתְּחִלָּה וָאִיקָץ׃}
}
{
	\eE{
		and they came inside, and it wasn’t known that they had
		come inside! And their appearance was bad, as at first.
		And I woke up.
	}
}
{
	Then they came inside and nobody knew and then he woke up.
}

\Verse{22}
{
	\hE{וָאֵרֶא בַּחֲלֹמִי וְהִנֵּה שֶׁבַע שִׁבֳּלִים עֹלֹת בְּקָנֶה אֶחָד מְלֵאֹת וְטֹבוֹת׃}
}
{
	\eE{
		“And I saw in my dream, and here were seven ears of grain
		coming up on one stalk, full and good.
	}
}
{
	Seven fluffy breadsticks.
}

\Verse{23}
{
	\hE{וְהִנֵּה שֶׁבַע שִׁבֳּלִים צְנֻמוֹת דַּקּוֹת שְׁדֻפוֹת קָדִים צֹמְחוֹת אַחֲרֵיהֶם׃}
}
{
	\eE{
		And here were seven ears of grain, dried up, thin,
		scorched by the east wind, growing after them.
	}
}
{
	Seven burned breadsticks.
}

\Verse{24}
{
	\hE{וַתִּבְלַעְןָ הַשִּׁבֳּלִים הַדַּקֹּת אֵת שֶׁבַע הַשִּׁבֳּלִים הַטֹּבוֹת וָאֹמַר אֶל הַחַרְטֻמִּים וְאֵין מַגִּיד לִי׃}
}
{
	\eE{
		And the thin ears of grain swallowed the seven good ears
		of grain. “And I said it to the magicians, and there’s no
		one who can tell me.”
	}
}
{
	Then the breadsticks ate the breadsticks and it was scary
	and the magicians are confused.
}

\Verse{25}
{
	\hE{וַיֹּאמֶר יוֹסֵף אֶל פַּרְעֹה חֲלוֹם פַּרְעֹה אֶחָד הוּא אֵת אֲשֶׁר הָאֱלֹהִים עֹשֶׂה הִגִּיד לְפַרְעֹה׃}
}
{
	\eE{
		And Joseph said to Pharaoh, “Pharaoh’s dream: it is one.
		It has told to Pharaoh what {\God} is doing.
	}
}
{
	So Joseph's like ``they're the same dream.''

	He gives them the old:

	\begin{enumerate}
	\item A = B
	\item Both A and B are God talking to you, not to me.
	\item Go figure it out you're the special one not me.\fB{
		Cf. Daniel 2:46-49.
	}
	\end{enumerate}

	

%	\footnote{
%		Pharaoh’s dream: it is one. It has told to Pharaoh what God is doing.
%		The pronoun is unclear. The verse could mean that God has told Pharaoh
%		what He is doing or that the dream has told Pharaoh what God is doing.
%		Even on the latter meaning, Joseph has already said that “meanings
%		belong to God” and that “God will answer regarding Pharaoh’s
%		well-being.” But the question is still whether God has provided both the
%		dream and the interpretation of it—or just the interpretation, while the
%		dream itself may derive from elsewhere. This question remains unanswered
%		in the text! (The same applies to 41:28, “It has shown Pharaoh what God
%		is doing.”)
%	}
}

\Verse{26}
{
	\hE{שֶׁבַע פָּרֹת הַטֹּבֹת שֶׁבַע שָׁנִים הֵנָּה וְשֶׁבַע הַשִּׁבֳּלִים הַטֹּבֹת שֶׁבַע שָׁנִים הֵנָּה חֲלוֹם אֶחָד הוּא׃}
}
{
	\eE{
		The seven good cows: they’re seven years. And the seven
		good ears of grain: they’re seven years. It’s one dream.
	}
}
{
	All the stuff is years.
	
	And cows are breadsticks.
}

\Verse{27}
{
	\hE{וְשֶׁבַע הַפָּרוֹת הָרַקּוֹת וְהָרָעֹת הָעֹלֹת אַחֲרֵיהֶן שֶׁבַע שָׁנִים הֵנָּה וְשֶׁבַע הַשִּׁבֳּלִים הָרֵקוֹת שְׁדֻפוֹת הַקָּדִים יִהְיוּ שֶׁבַע שְׁנֵי רָעָב׃}
}
{
	\eE{
		And the seven cows that were scrawny and bad that were
		coming up after them: they’re seven years. And the seven
		ears of grain that were scrawny, scorched by the east
		wind, will be seven years of famine.
	}
}
{
	Bad cows are bad years.

	Bad breadsticks are bad years.

	Food with problems means problems with food.
}

\Verse{28}
{
	\hE{הוּא הַדָּבָר אֲשֶׁר דִּבַּרְתִּי אֶל פַּרְעֹה אֲשֶׁר הָאֱלֹהִים עֹשֶׂה הֶרְאָה אֶת פַּרְעֹה׃}
}
{
	\eE{
		That is the thing that I spoke to Pharaoh: it has shown
		Pharaoh what {\God} is doing.
	}
}
{
	That's the deal.
	
	That's what {\God}'s up to.
}

\Verse{29}
{
	\hE{הִנֵּה שֶׁבַע שָׁנִים בָּאוֹת שָׂבָע גָּדוֹל בְּכל אֶרֶץ מִצְרָיִם׃}
}
{
	\eE{
		Here, seven years are coming: big bountifulness in all the
		land of Egypt.
	}
}
{
	In summary, the next seven years you'll have food with no problems,
	which as we saw above means no problems with food.
}

\Verse{30}
{
	\hE{וְקָמוּ שֶׁבַע שְׁנֵי רָעָב אַחֲרֵיהֶן וְנִשְׁכַּח כּל הַשָּׂבָע בְּאֶרֶץ מִצְרָיִם וְכִלָּה הָרָעָב אֶת הָאָרֶץ׃}
}
{
	\eE{
		And seven years of famine will rise after them, and all
		the bountifulness in the land of Egypt will be forgotten,
		and the famine will finish off the land.
	}
}
{
	After that, seven years of problem food.
}

\Verse{31}
{
	\hE{וְלֹא יִוָּדַע הַשָּׂבָע בָּאָרֶץ מִפְּנֵי הָרָעָב הַהוּא אַחֲרֵי כֵן כִּי כָבֵד הוּא מְאֹד׃}
}
{
	\eE{
		And the bountifulness won’t be known in the land on
		account of that famine afterwards because it will be very
		heavy.
	}
}
{
	People will forget about the good times because
	those cows are in the past and they're hungry
	now because bread problems, any questions?
}

\Verse{32}
{
	\hE{וְעַל הִשָּׁנוֹת הַחֲלוֹם אֶל פַּרְעֹה פַּעֲמָיִם כִּי נָכוֹן הַדָּבָר מֵעִם הָאֱלֹהִים וּמְמַהֵר הָאֱלֹהִים לַעֲשֹׂתוֹ׃}
}
{
	\eE{
		And about the recurrence of the dream to Pharaoh two
		times: because the thing is right from {\God}, and {\God} is
		hurrying to do it.
	}
}
{
	You dreaming it twice means it's coming soon.
	
	Number equals velocity in dream space.
%	\footnote{
%		right from God. This term for “right” (Hebrew nakôn) will be recalled to
%		play an important part in the story of the exodus (Exod 8:22; see the
%		text and comment there).
%	}
}

\Verse{33}
{
	\hE{וְעַתָּה יֵרֶא פַרְעֹה אִישׁ נָבוֹן וְחָכָם וִישִׁיתֵהוּ עַל אֶרֶץ מִצְרָיִם׃}
}
{
	\eE{
		And now, let Pharaoh look out for an understanding and
		wise man and set him over the land of Egypt.
	}
}
{
	You should really find a wise guy to help out with all that.
}

\Verse{34}
{
	\hE{יַעֲשֶׂה פַרְעֹה וְיַפְקֵד פְּקִדִים עַל הָאָרֶץ וְחִמֵּשׁ אֶת אֶרֶץ מִצְרַיִם בְּשֶׁבַע שְׁנֵי הַשָּׂבָע׃}
}
{
	\eE{
		Let Pharaoh do it and appoint overseers over the land and
		five-out the land of Egypt in the seven years of
		bountifulness.
	}
}
{
	You should pick a team to plan ahead
	during the first seven good cows.

	Have them five-out\fC{No one knows what this means.} the land.

%	\footnote{
%		five-out. This perhaps means to divide the land into portions of five
%		and save one-fifth for the bad years, but the meaning is uncertain
%		because this is the only occurrence of this verb in the Hebrew Bible.
%	}
}

\Verse{35}
{
	\hE{וְיִקְבְּצוּ אֶת כּל אֹכֶל הַשָּׁנִים הַטֹּבוֹת הַבָּאֹת הָאֵלֶּה וְיִצְבְּרוּ בָר תַּחַת יַד פַּרְעֹה אֹכֶל בֶּעָרִים וְשָׁמָרוּ׃}
}
{
	\eE{
		And let them gather all the food of these coming good
		years and pile up grain under Pharaoh’s hand, food in the
		cities, and watch over it.
	}
}
{
	Should probably save some cows for when the cows come.
}

\Verse{36}
{
	\hE{וְהָיָה הָאֹכֶל לְפִקָּדוֹן לָאָרֶץ לְשֶׁבַע שְׁנֵי הָרָעָב אֲשֶׁר תִּהְיֶיןָ בְּאֶרֶץ מִצְרָיִם וְלֹא תִכָּרֵת הָאָרֶץ בָּרָעָב׃}
}
{
	\eE{
		And the food will become a deposit for the land for the
		seven years of famine that will be in the land of Egypt so
		the land won’t be cut off in the famine.”
	}
}
{
	Like a savings accownt.
}

\Verse{37}
{
	\hE{וַיִּיטַב הַדָּבָר בְּעֵינֵי פַרְעֹה וּבְעֵינֵי כּל עֲבָדָיו׃}
}
{
	\eE{
		And the thing was good in Pharaoh’s eyes and in all his
		servants’ eyes.
	}
}
{
	The king agreed.
}

\Verse{38}
{
	\hE{וַיֹּאמֶר פַּרְעֹה אֶל עֲבָדָיו הֲנִמְצָא כָזֶה אִישׁ אֲשֶׁר רוּחַ אֱלֹהִים בּוֹ׃}
}
{
	\eE{
		And Pharaoh said to his servants, “Will we find such a man
		as this, that {\God}’s spirit is in him?”
	}
}
{
	Pharaoh looks around \aB{\textsl{like an absolute cartoon}} going
	
	``Where are we gonna find a guy like \textsl{that}?''

	Starting to see why Joseph had such an easy time getting
	put in charge of everything.
}

\Verse{39}
{
	\hE{וַיֹּאמֶר פַּרְעֹה אֶל יוֹסֵף אַחֲרֵי הוֹדִיעַ אֱלֹהִים אוֹתְךָ אֶת כּל זֹאת אֵין נָבוֹן וְחָכָם כָּמוֹךָ׃}
}
{
	\eE{
		And Pharaoh said to Joseph, “Since {\God} has made you know
		all this, there’s no one as understanding and wise as you.
	}
}
{
	Then Pharaoh's like,

	``I have an idea...
	
	Why don't \textsc{You} be the guy!''

	\aB{I picture all the advisors giving some gentle applause here.}

%	\footnote{
%		understanding and wise. Some interpreters identify the Joseph story as a
%		wisdom tale. It is treated as comparable to wisdom tales from numerous
%		cultures about a wise young man who is successful in a royal court
%		because of his cleverness. The story of Mordecai in the book of Esther
%		and the story of Daniel are compared in this context as well. I would
%		caution that there is a danger here of overdefining the word “wisdom.”
%		The term was commonly used to refer to the books of Proverbs, Job, and
%		Ecclesiastes, works similar to philosophy; composed, like the
%		pre-Socratics, in poetry; sometimes set in very loose story contexts
%		(Job, Ecclesiastes), but basically dealing with issues outright, not
%		developing them through a prose narrative fabric. Now, in the first
%		place, if we redefine wisdom to include the Joseph cycle, and thus the
%		books of Daniel and Esther, we shall be forced to include a number of
%		other comparable sections of the Hebrew Bible under this heading as
%		well. We shall be left with a category of works with more elements of
%		variety than commonality (Ecclesiastes and Esther?), and the term
%		“wisdom” will be rendered less descriptive and useful. Second, the facts
%		of the Joseph story do not really fit this model in any case. Joseph is
%		not successful initially because of his cleverness or his sage advice.
%		He succeeds because he can interpret dreams, apparently by revelation
%		rather than by his own insight, since even the most brilliant Freudian
%		analyst would be unlikely to guess that a dream of cows eating cows
%		means that there will be prosperity followed by famine in the Near East.
%		Joseph is described as stating that it is God and not he who directs
%		events and sheds light on dreams (40:8; 41:16,25; 50:20). One might
%		suggest that this story is still based on the wisdom tale model but has
%		been adapted to fit the biblical author’s literary and theological
%		requirements. That would be an interesting phenomenon to pursue, but it
%		would be a study of adaptation, not of wisdom. To understand biblical
%		wisdom, we are better advised to study wisdom literature as it has
%		classically been understood: Proverbs, Job, and Ecclesiastes.
%	}
}

\Verse{40}
{
	\hE{אַתָּה תִּהְיֶה עַל בֵּיתִי וְעַל פִּיךָ יִשַּׁק כּל עַמִּי רַק הַכִּסֵּא אֶגְדַּל מִמֶּךָּ׃}
}
{
	\eE{
		You shall be over my house, and at your mouth all my
		people shall conform. I shall be greater than you only in
		the throne.”
	}
}
{
	Yeah. And I'll put you in charge of my house!

	And everybody will do what you say!

	I'll only be higher status as sort of a formal thing.\fB{
		He. Did. Not. Say. This.
	}

	Like a figurehead!\fB{
		\textsl{This} is Ala 1:05:22.
	}
	\fC{
		REF: conform. The verb here, which usually means “to kiss,”
		is a problem for all translators. It may reflect a scribal 
		problem or an idiom that is no longer familiar to us.
		\aB{
			\footnotesize Or the magic stick is making him say it.
		}
	}
}

\Verse{41}
{
	\hE{וַיֹּאמֶר פַּרְעֹה אֶל יוֹסֵף רְאֵה נָתַתִּי אֹתְךָ עַל כּל אֶרֶץ מִצְרָיִם׃}
}
{
	\eE{
		And Pharaoh said to Joseph, “See, I’ve put you over all
		the land of Egypt.”
	}
}
{
	The vizier now runs the kingdom.
}

\Verse{42}
{
	\hE{וַיָּסַר פַּרְעֹה אֶת טַבַּעְתּוֹ מֵעַל יָדוֹ וַיִּתֵּן אֹתָהּ עַל יַד יוֹסֵף וַיַּלְבֵּשׁ אֹתוֹ בִּגְדֵי שֵׁשׁ וַיָּשֶׂם רְבִד הַזָּהָב עַל צַוָּארוֹ׃}
}
{
	\eE{
		And Pharaoh took off his ring from his hand and put it on
		Joseph’s hand, and he had him dressed in linen garments,
		and he set a gold chain on his neck,
	}
}
{
	The Pharaoh gives him the ring off his hand,
	and some fancy royal clothes and some gold.

	\aB{Remember, all this is because he said cows are years.}

}

\Verse{43}
{
	\hE{וַיַּרְכֵּב אֹתוֹ בְּמִרְכֶּבֶת הַמִּשְׁנֶה אֲשֶׁר לוֹ וַיִּקְרְאוּ לְפָנָיו אַבְרֵךְ וְנָתוֹן אֹתוֹ עַל כּל אֶרֶץ מִצְרָיִם׃}
}
{
	\eE{
		and he had him driven in the chariot of the
		second-in-command whom he had, and they called in front of
		him, “Kneel,” so putting him over all the land of Egypt.
	}
}
{
	Then the king gives him the vice-president's car.

	And he drives around the kingdom with guys shouting
	``Kneel'' at people on the street.%
	\fA{
		The word translated ``Kneel'' here is \heb{אברך}.
		Nobody knows what it means, but it is often
		translated as ``kneel,'' due to 1. context
		2. some Rabbinic interpretations, and
		3. similarity to the word \heb{ברך}, which means
		both ``to kneel'' and ``to bless.''
	}
	\fB{It's the word from that ``I won't let you go until
	you \heb{ברך} me'' scene, spoken by Jacob (or maybe his
	unidentified night wrestling partner) right at the
	end before he (who?) \heb{ברך}s him. Welcome to the bible.
	It's wild in here.}
	\fC{Despite the common claim that \textit{abrek} is Egyptian for 
	`kneel,'' the required Egyptian verb *brk `kneel'' is unattested,
	and the initial \textit{a-} is morphologically suspect. The 
	resemblance is actually stronger within Semitic: Hebrew itself has 
	BRK `kneel/make kneel,'' while Syriac \textit{ʾabrek} and Arabic%
	\textit{ʾabraka} mean `make kneel.'' The term may simply be a Hebrew 
	speaking author dressing up a Semitic root to sound foreign.
	}


%	\footnote{
%		Kneel. Hebrew ‘abrk. Its meaning is unknown. Some see a derivation from
%		an Egyptian term. One rabbinic interpretation takes it to mean “Bend the knee.”
%	}
}

\Verse{44}
{
	\hE{וַיֹּאמֶר פַּרְעֹה אֶל יוֹסֵף אֲנִי פַרְעֹה וּבִלְעָדֶיךָ לֹא יָרִים אִישׁ אֶת יָדוֹ וְאֶת רַגְלוֹ בְּכל אֶרֶץ מִצְרָיִם׃}
}
{
	\eE{
		And Pharaoh said to Joseph, “I am Pharaoh, and without you
		a man won’t lift his hand and his foot in all the land of
		Egypt.”
	}
}
{
	Pharaoh's like ``I am Pharaoh.''

	Then he says nobody's allowed to walk or do hand stuff
	anymore without Joseph.
}

\Verse{45}
{
	\hE{וַיִּקְרָא פַרְעֹה שֵׁם יוֹסֵף צָפְנַת פַּעְנֵחַ וַיִּתֶּן לוֹ אֶת אָסְנַת בַּת פּוֹטִי פֶרַע כֹּהֵן אֹן לְאִשָּׁה וַיֵּצֵא יוֹסֵף עַל אֶרֶץ מִצְרָיִם׃}
}
{
	\eE{
		And Pharaoh called Joseph’s name Zaphenathpaneah. And he
		gave him Asenath, daughter of Potiphera, priest of On, as
		a wife. And Joseph went out over the land of Egypt.
	}
}
{
	\aB{Pharaoh renames Joseph to some gibberish.}

	So the king's like:
	
	Welcome to Egypt, your name will no longer be

	\egypt{𓇋𓅱𓊃𓆑𓀀}

	but rather

	\egypt{𓋴𓆑𓈖𓏏𓀜 𓊪𓂝𓈖𓐍𓀭𓅭𓇳}.


	\aB{
		Perhaps why this renaming event is less well-remembered
		than the Abraham and Jacob ones.\fB{
			Isaac didn't need a new name,
			his name was already funny.
		}
	}

	He also gives him a free wife.

	Joseph goes out to explore this new country he sort of conquered.

	\aB{To reiterate, all this is because he said cows are years.}

%	\footnote{
%		Asenath, daughter of Poti-phera, priest of On. Like Moses, Joseph
%		marries the daughter of a priest. The Pharaoh would naturally provide
%		Joseph with a prestigious marriage, but he would be disinclined to let
%		Joseph, a foreigner, marry anyone from the royal house. Marriage to a
%		woman from a distinguished priestly family would be the next best thing.
%	}
}

\Verse{46}
{
	\hP{וְיוֹסֵף בֶּן שְׁלֹשִׁים שָׁנָה בְּעמְדוֹ לִפְנֵי פַּרְעֹה מֶלֶךְ מִצְרָיִם}
	\hR{וַיֵּצֵא יוֹסֵף מִלִּפְנֵי פַרְעֹה וַיַּעֲבֹר בְּכל אֶרֶץ מִצְרָיִם׃}
}
{
	\eP{And Joseph was thirty years old when he stood in front of Pharaoh, king of Egypt.}
	\eR{And Joseph went out from in front of Pharaoh and passed through all the land of Egypt.}
}
{
}

\Verse{47}
{
	\hE{וַתַּעַשׂ הָאָרֶץ בְּשֶׁבַע שְׁנֵי הַשָּׂבָע לִקְמָצִים׃}
}
{
	\eE{
		And the land produced in the seven years of bountifulness
		by fistfuls.
	}
}
{
}

\Verse{48}
{
	\hE{וַיִּקְבֹּץ אֶת כּל אֹכֶל שֶׁבַע שָׁנִים אֲשֶׁר הָיוּ בְּאֶרֶץ מִצְרַיִם וַיִּתֶּן אֹכֶל בֶּעָרִים אֹכֶל שְׂדֵה הָעִיר אֲשֶׁר סְבִיבֹתֶיהָ נָתַן בְּתוֹכָהּ׃}
}
{
	\eE{
		And he gathered all the food of the seven years, which
		were in the land of Egypt, and he put food in cities; he
		put the food of the city’s fields that were around it
		inside of it.
	}
}
{
}

\Verse{49}
{
	\hE{וַיִּצְבֹּר יוֹסֵף בָּר כְּחוֹל הַיָּם הַרְבֵּה מְאֹד עַד כִּי חָדַל לִסְפֹּר כִּי אֵין מִסְפָּר׃}
}
{
	\eE{
		And Joseph piled up grain like the sand of the sea, very
		much, until he stopped counting, because it was without
		number.
	}
}
{
}

\Verse{50}
{
	\hE{וּלְיוֹסֵף יֻלַּד שְׁנֵי בָנִים בְּטֶרֶם תָּבוֹא שְׁנַת הָרָעָב אֲשֶׁר יָלְדָה לּוֹ אָסְנַת בַּת פּוֹטִי פֶרַע כֹּהֵן אוֹן׃}
}
{
	\eE{
		And two sons were born to Joseph before the year of famine
		came—to whom Asenath, daughter of Poti-phera, priest of
		On, gave birth for him.
	}
}
{
}

\Verse{51}
{
	\hE{וַיִּקְרָא יוֹסֵף אֶת שֵׁם הַבְּכוֹר מְנַשֶּׁה כִּי נַשַּׁנִי אֱלֹהִים אֶת כּל עֲמָלִי וְאֵת כּל בֵּית אָבִי׃}
}
{
	\eE{
		And Joseph called the firstborn’s name Manasseh, “because
		{\God} has made me forget all of my trouble and all of my
		father’s house.”
	}
}
{
}

\Verse{52}
{
	\hE{וְאֵת שֵׁם הַשֵּׁנִי קָרָא אֶפְרָיִם כִּי הִפְרַנִי אֱלֹהִים בְּאֶרֶץ ענְיִי׃}
}
{
	\eE{
		And he called the second one’s name Ephraim, “because {\God}
		has made me fruitful in the land of my degradation.”
	}
}
{
}

\Verse{53}
{
	\hE{וַתִּכְלֶינָה שֶׁבַע שְׁנֵי הַשָּׂבָע אֲשֶׁר הָיָה בְּאֶרֶץ מִצְרָיִם׃}
}
{
	\eE{
		And the seven years of the bountifulness that was in the
		land of Egypt finished,
	}
}
{
}

\Verse{54}
{
	\hE{וַתְּחִלֶּינָה שֶׁבַע שְׁנֵי הָרָעָב לָבוֹא כַּאֲשֶׁר אָמַר יוֹסֵף וַיְהִי רָעָב בְּכל הָאֲרָצוֹת וּבְכל אֶרֶץ מִצְרַיִם הָיָה לָחֶם׃}
}
{
	\eE{
		and the seven years of famine started to come, as Joseph
		had said. And there was famine in all the lands, and in
		all the land of Egypt there was bread.
	}
}
{
}

\Verse{55}
{
	\hE{וַתִּרְעַב כּל אֶרֶץ מִצְרַיִם וַיִּצְעַק הָעָם אֶל פַּרְעֹה לַלָּחֶם וַיֹּאמֶר פַּרְעֹה לְכל מִצְרַיִם לְכוּ אֶל יוֹסֵף אֲשֶׁר יֹאמַר לָכֶם תַּעֲשׂוּ׃}
}
{
	\eE{
		And all the land of Egypt hungered, and the people cried
		to Pharaoh for the bread, and Pharaoh said to all Egypt,
		“Go to Joseph. Do what he’ll tell you.”
	}
}
{
}

\Verse{56}
{
	\hE{וְהָרָעָב הָיָה עַל כּל פְּנֵי הָאָרֶץ וַיִּפְתַּח יוֹסֵף אֶת כּל אֲשֶׁר בָּהֶם וַיִּשְׁבֹּר לְמִצְרַיִם וַיֶּחֱזַק הָרָעָב בְּאֶרֶץ מִצְרָיִם׃}
}
{
	\eE{
		And the famine was on all the face of the earth. And
		Joseph opened everything that was in them and sold to
		Egypt. And the famine was strong in the land of Egypt.
	}
}
{
%	\footnote{
%		everything that was in them. In what? The Septuagint reads, “Joseph
%		opened all the grain storehouses” (), which is consistent with 41:49.
%	}
}

\Verse{57}
{
	\hE{וְכל הָאָרֶץ בָּאוּ מִצְרַיְמָה לִשְׁבֹּר אֶל יוֹסֵף כִּי חָזַק הָרָעָב בְּכל הָאָרֶץ׃}
}
{
	\eE{
		And all the earth came to Egypt to buy, to Joseph, because
		the famine was strong in all the earth.
	}
}
{
}
